\documentclass[11pt]{article}
\usepackage[margin=1in]{geometry}
\usepackage{amsmath,amssymb}
\usepackage{xcolor} % \textcolor{red}{ADDED COLOR SUPPORT FOR THE NEW SUMMARY.}

\title{Literature Review Summaries}
\author{}
\date{}

\begin{document}
\maketitle

% Labels are maintained in ../lit_review/CITATION_LABELS.md.
\section{[HMR] Quantile-based Loss Filtering\\for Outlier-Robust Stochastic Gradient Descent}
\label{sec:hmr}
\textbf{Authors:} Jamie Haddock, Anna Ma, and Elizaveta Rebrova.

The paper studies finite-sum optimization with corrupted component functions,
aiming to recover a minimizer of the uncorrupted objective. Its main
contribution is Quantile-$k$-Loss SGD (Q$k$L-SGD): each iteration samples $k$
component losses, selects an index uniformly from their lower empirical
$q$-quantile, and takes one gradient step. The method interpolates between
minimum-$k$-loss SGD and ordinary SGD, balancing outlier rejection against
selection of informative updates.

The authors establish two complementary convergence results. For sufficiently
large samples relative to the number of corruptions, they prove linear
convergence in expected squared distance to the clean minimizer under
smoothness, interpolation, subset strong convexity, and suitable sampling and
step-size conditions. For arbitrary sample sizes, a probabilistic analysis
relates convergence to the chance of selecting an outlier and the expected
curvature of a good update. With bounded corrupted updates and sufficient
contraction, this yields linear convergence to an error neighborhood; exact
convergence requires additional control of corrupted updates near the solution.

Experiments on polynomial regression, regularized logistic regression, and
regularized hinge loss show that intermediate quantiles often outperform
both ordinary SGD and minimum-$k$-loss SGD. The latter can stall by repeatedly
selecting nearly solved components. The paper identifies averaging accepted
gradients, momentum, and variance reduction as directions for future work.


\section{[SR] Beyond the error horizon: averaged quantile Kaczmarz methods for corrupted and noisy linear systems}
\label{sec:sr}
\textbf{Authors:} Sofiia Shvaiko and Elizaveta Rebrova.

The paper studies linear systems whose right-hand sides contain both sparse,
potentially large corruptions and dense noise. Its method, averaged quantile
randomized Kaczmarz (AQRK), selects equations with small absolute residuals
and averages a batch of relaxed single-row updates. A central observation is
that filtering can bias the accepted noise: averaging reduces fluctuations
but cannot by itself remove this bias. When repeated observations are
available, AQRK queries the selected equations again, keeping the update
noise independent of the selection event while leaving the sparse
corruptions fixed.

For full screening, independent bounded mean-zero query noise, a sufficiently
small corruption fraction relative to the matrix's restricted conditioning,
and corruptions sufficiently large relative to the noise, the authors prove
\[
\mathbb E\|x_k-x^\star\|^2
\le (1-\eta\alpha)^k\|x_0-x^\star\|^2
+\frac{\alpha r_\infty^2}{\eta s},
\]
where $\alpha$ is an admissible stepsize, $s$ is the update batch size,
$r_\infty$ bounds the query noise, and $\eta>0$ depends on the matrix,
quantile, and corruption fraction. Thus the root-mean-square error horizon
scales as $\sqrt{\alpha/s}$. Suitable decreasing-stepsize or increasing-batch
schedules yield mean-square convergence to the clean solution. For
Gaussian-like designs, the horizon is of order
$\sqrt{\alpha n/s}\,r_\infty$, so a batch of order $n$ with an admissible
constant stepsize reduces it to the noise scale.

Experiments on synthetic systems support the batch-size scaling and examine
the importance of corruption separation. Tomography tests show improvements
on an ill-conditioned system beyond the theoretical assumptions, including
fixed noise, where filtering bias can persist and semi-convergence appears.
As an inner solver for CP tensor fitting, AQRK also recovers corrupted
low-rank tensors in cases where standard alternating least squares,
unfiltered averaging, and single-row quantile updates fail. Future directions
include fixed-noise analysis, weaker separation conditions, optimal parameter
schedules, and stopping rules for changing least-squares subproblems.
\par


\end{document}
